% Compile with pdfLaTeX twice.
\documentclass[a4paper,11pt]{amsart}
\usepackage{amssymb,mathtools}
\usepackage{microtype}
\usepackage{xcolor}
\definecolor{citeblue}{RGB}{0,0,128}
\definecolor{refred}{RGB}{128,0,0}
\usepackage[colorlinks=true,citecolor=citeblue,linkcolor=refred,urlcolor=citeblue]{hyperref}
\setlength{\headheight}{15.2pt}
\setlength{\emergencystretch}{1em}
\setlength{\textwidth}{15cm}
\setlength{\textheight}{22.75cm}
\setlength{\topmargin}{0.5cm}
\setlength{\oddsidemargin}{0.5cm}
\setlength{\evensidemargin}{0.5cm}
\pagestyle{plain}
\renewcommand{\bibliofont}{\small}
\hypersetup{pdftitle={Standard graded nearly Gorenstein Ehrhart rings have unimodal h-vectors},pdfauthor={Takayuki Hibi and Sora Miyashita},pdfsubject={Canonical trace ideals, Frobenius splitting, Ehrhart rings, and h-vectors},pdfkeywords={Nearly Gorenstein ring, Ehrhart ring, canonical trace ideal, Frobenius splitting, strong Lefschetz property, graded domain, unimodality, level ring}}

\DeclareMathOperator{\tr}{tr}
\DeclareMathOperator{\conv}{conv}
\DeclareMathOperator{\Hilb}{Hilb}
\DeclareMathOperator{\codeg}{codeg}
\DeclareMathOperator{\indeg}{indeg}
\DeclareMathOperator{\relint}{relint}
\DeclareMathOperator{\Hom}{Hom}
\DeclareMathOperator{\height}{ht}
\DeclareMathOperator{\depth}{depth}
\newcommand{\kk}{k}
\newcommand{\QQ}{\mathbb{Q}}
\newcommand{\ZZ}{\mathbb{Z}}
\newcommand{\RR}{\mathbb{R}}
\newcommand{\m}{\mathfrak{m}}
\newcommand{\natur}{\natural}

\newtheorem{theorem}{Theorem}[section]
\newtheorem{lemma}[theorem]{Lemma}
\newtheorem{proposition}[theorem]{Proposition}
\newtheorem{corollary}[theorem]{Corollary}
\newtheorem{conjecture}[theorem]{Conjecture}
\newtheorem{question}[theorem]{Question}
\newtheorem{maintheorem}{Theorem}
\renewcommand{\themaintheorem}{\Alph{maintheorem}}
\theoremstyle{definition}
\newtheorem{example}[theorem]{Example}
\newtheorem{remark}[theorem]{Remark}

\title[Unimodality for nearly Gorenstein Ehrhart rings]{Standard graded nearly Gorenstein Ehrhart rings have unimodal $h$-vectors}
\author{Takayuki Hibi}
\author{Sora Miyashita}
\address{Department of Pure and Applied Mathematics, Graduate School of Information Science and Technology, The University of Osaka, Suita, Osaka 565-0871, Japan}
\email[T.H.]{hibi@math.sci.osaka-u.ac.jp}
\email[S.M.]{u804642k@ecs.osaka-u.ac.jp}
\date{November 21, 2026}
\subjclass[2020]{Primary 13H10, 05E40; Secondary 13A02, 52B20}
\keywords{Nearly Gorenstein ring, Ehrhart ring, canonical trace ideal, Frobenius splitting, strong Lefschetz property, graded domain, unimodality, level ring}

\begin{document}

\begin{abstract}
We prove that the $h$-vector $(h_0,\ldots,h_s)$ of every standard graded nearly Gorenstein Ehrhart ring is unimodal, with a peak at position $\lceil s/2\rceil$. More generally, a standard graded Ehrhart ring has an alternatingly increasing $h$-vector if the radical of the ideal generated by the degree-one part of its canonical trace ideal contains the homogeneous maximal ideal. In characteristic zero or two, the generic Artinian reduction of each such ring has the strong Lefschetz property. We also prove unimodality for every standard graded nearly Gorenstein domain of dimension at least two with $h$-vector $(1,h_1,\ldots,h_s)$, where $s\geq2$, provided that $h_1-h_s\leq1$. In the level case this $h$-vector is both concave and alternatingly increasing. If $h_1=h_s$, then $h_1=\cdots=h_s$.
\end{abstract}

\maketitle

\section{Introduction}

Let $R=\bigoplus_{j\geq0} R_j$ be a standard graded Cohen--Macaulay algebra over a field $\kk$. Write $d=\dim R$ and let $\m_R=\bigoplus_{j>0} R_j$ be its \emph{homogeneous maximal ideal}. The \emph{Hilbert series} $\Hilb_R(t)=\sum_{j\geq0}(\dim_{\kk} R_j)t^j$ has the form $h_R(t)/(1-t)^d$, where $h_R(t)=1+h_1t+\cdots+h_st^s$ and $h_s\neq0$. We call $h_R(t)$ the \emph{$h$-polynomial} and $h(R)=(1,h_1,\ldots,h_s)$ the \emph{$h$-vector} of $R$. The $h$-vector is \emph{unimodal} if $h_0\leq\cdots\leq h_j\geq\cdots\geq h_s$ for some $j$ with $0\leq j\leq s$. It is \emph{log-concave} if $h_i^2\geq h_{i-1}h_{i+1}$ for every $i$ with $1\leq i<s$.

By classical results of Stanley, the $h$-vector of every standard graded Gorenstein algebra is symmetric \cite[Theorem~4.1]{S78}, and every standard graded Cohen--Macaulay domain with a symmetric $h$-vector is Gorenstein \cite[Theorem~4.4]{S78}. A symmetric $h$-vector need not be unimodal. For example, Stanley constructed an Artinian Gorenstein algebra with $h$-vector $(1,13,12,13,1)$ \cite[Example~4.3]{S78}.

The following conjecture goes back to Hibi and Stanley and remains open.

\begin{conjecture}[{\cite[Conjecture~1.5]{H89}, \cite[Conjecture~4~(a), p.~512]{S89}}]\label{conj:classical}
Every standard graded Gorenstein domain over a field has a unimodal $h$-vector.
\end{conjecture}

Stanley's original conjecture asked for the stronger property of log-concavity for standard graded Cohen--Macaulay domains, with the Gorenstein case suggested as a possible restriction. Conca, Singh, and Varbaro constructed semi-standard graded Gorenstein normal domains with non-unimodal $h$-vectors \cite[Theorem~4.2]{CSV25}. Their examples show that the standard grading hypothesis in Conjecture~\ref{conj:classical} is essential.

Let $P\subset\RR^q$ be a lattice polytope. Its \emph{Ehrhart ring} $A_P=\kk[\mathbf{x}^{\boldsymbol u}t^n:\boldsymbol u\in nP\cap\ZZ^q,\ n\geq0]$, graded by the exponent of $t$, is a normal Cohen--Macaulay domain \cite[Theorem~1]{Hoc72}. It is standard graded precisely when $P$ has the \emph{integer decomposition property} (IDP): for every $n\geq1$, each lattice point of $nP$ is a sum of $n$ lattice points of $P$. The Hilbert series of $A_P$ has the form $h_P^*(t)/(1-t)^{\dim P+1}$. The numerator is the \emph{$h^*$-polynomial} $h_P^*(t)=\sum_{i=0}^s h_it^i$, where $h_s\ne0$. Its coefficient vector $h^*(P)=(h_0,\ldots,h_s)$ is the \emph{$h^*$-vector} of $P$. If $P$ is IDP, then $h^*(P)=h(A_P)$. We write $\relint$ for the \emph{relative interior} and take interiors in the affine span of $P$. We call $P$ \emph{Gorenstein} if $A_P$ is Gorenstein.

In 1992 Hibi conjectured unimodality of $h^*(P)$ for every reflexive lattice polytope \cite[p.~111]{H92}. Here a full-dimensional lattice polytope with the origin in its interior is \emph{reflexive} if its polar is a lattice polytope. Musta\c{t}\u{a} and Payne gave a counterexample \cite[Example~1.1]{MP05}, and Payne constructed further counterexamples among reflexive simplices \cite[Theorem~1.4]{P08}. These examples fail IDP. For standard graded Gorenstein toric rings, unimodality was also conjectured by Ohsugi and Hibi \cite[Introduction]{OH06}.

Bruns and R\"omer proved unimodality for Gorenstein polytopes with a regular unimodular triangulation \cite[Theorem~1]{BR07}. Adiprasito, Papadakis, Petrotou, and Steinmeyer announced results removing the triangulation hypothesis for Gorenstein IDP polytopes \cite{APPS22}. Adiprasito, Papadakis, and Petrotou proved unimodality for reflexive IDP polytopes \cite[Theorem~1.4]{APP25} (see also \cite[Corollary~6.3]{AKOPP26}). By \cite[Corollary~7]{BR07}, every Gorenstein IDP polytope has the same $h^*$-polynomial as a reflexive IDP polytope, so unimodality holds for every Gorenstein IDP polytope. On the other hand, Ferroni constructed smooth IDP polytopes with non-unimodal $h^*$-vectors and Gorenstein IDP polytopes with non-log-concave $h^*$-vectors \cite{F26}.

This raises the question whether the $h$-vector remains unimodal when a standard graded Ehrhart ring is nearly Gorenstein. Let $\omega_R$ denote the \emph{graded canonical module} of $R$. The \emph{canonical trace ideal} is $\tr_R(\omega_R)=\sum_{\varphi\in\Hom_R(\omega_R,R)}\varphi(\omega_R)$, where the sum runs over homogeneous maps of all degrees. The ring $R$ is Gorenstein if and only if $\tr_R(\omega_R)=R$ \cite[Lemma~2.1]{HHS19}. A Cohen--Macaulay ring $R$ is \emph{nearly Gorenstein} if $\tr_R(\omega_R)\supseteq\m_R$ \cite{HHS19}. The ring $R$ is \emph{level} if $\omega_R$ is generated in one degree and \emph{anticanonical level} if $\Hom_R(\omega_R,R)$ is generated in one degree (cf.\ \cite[Definition~3.3]{Pa19}).

For IDP $(0,1)$-polytopes $P$, the nearly Gorenstein property of $A_P$ was characterized in \cite[Theorem~34]{HKMM23}.

We call an $h$-vector \emph{alternatingly increasing} if its entries are nondecreasing when read alternately from the two ends, starting on the left.

We use the following condition on the canonical trace ideal \cite[Definition~1.1]{M26}:
\[
(\natur)\qquad \sqrt{[\tr_R(\omega_R)]_1R}\supseteq\m_R.
\]
The following is our first main result.

\begin{maintheorem}\label{thm:A}
Let $P$ be an IDP lattice polytope and let $R=A_P$ have $h$-vector $h(R)=(h_0,\ldots,h_s)$. Suppose that one of the following conditions holds:
\begin{enumerate}
\item $R$ satisfies $(\natur)$;
\item $R$ is level and anticanonical level, and $[\tr_R(\omega_R)]_1\neq0$;
\item $R$ is level and $P$ has an interior lattice point.
\end{enumerate}
Then $h(R)$ is alternatingly increasing, that is,
\[
 h_0\leq h_s\leq h_1\leq h_{s-1}\leq\cdots\leq h_{\lceil s/2\rceil}.
\]
In particular, $h_0\leq\cdots\leq h_{\lceil s/2\rceil}\geq\cdots\geq h_s$, so the $h$-vector is unimodal, with a peak at position $\lceil s/2\rceil$.
\end{maintheorem}

If $R$ is nearly Gorenstein, then $[\tr_R(\omega_R)]_1=R_1$. Thus $(\natur)$ holds, and Theorem~\ref{thm:A}~\textup{(1)} gives the following corollary.

\begin{corollary}\label{cor:ng-ehrhart}
The $h$-vector $(h_0,\ldots,h_s)$ of a standard graded nearly Gorenstein Ehrhart ring is unimodal, with a peak at position $\lceil s/2\rceil$.
\end{corollary}

In characteristic zero or two, the generic Artinian reduction of $R$ in Theorem~\ref{thm:A} has the strong Lefschetz property (Corollary~\ref{cor:ehrhart-consequences}\,\textup{(\ref{item:ehrhart-lefschetz})}). In the Gorenstein case this follows from \cite[Theorem~6.1]{APP25} and the openness statement in \cite[Proposition~13.6]{PP20}.

For level Ehrhart rings, Theorem~\ref{thm:A}~\textup{(3)} answers the unimodality question in \cite{Hi19} and also follows from \cite{APP25} (see Remark~\ref{rem:app-scope}).

For a standard graded nearly Gorenstein domain of dimension $d\ge2$, $h_s=1$ forces $R$ to be Gorenstein \cite[Theorem~1.2]{MPG26}. We compare $h_s$ with the \emph{codimension} $h_1=\dim_{\kk}R_1-d$. An $h$-vector is \emph{concave} if $h_i-h_{i-1}\geq h_{i+1}-h_i$ for every $1\leq i<s$.

The following is our second main result.

\begin{maintheorem}\label{thm:B}
Let $R$ be a standard graded nearly Gorenstein domain over a field, with $\dim R\geq2$ and $h$-vector $h(R)=(1,h_1,\ldots,h_s)$, where $s\geq2$. Assume that $h_1-h_s\leq1$. Then $h_s\leq h_1$, and $h(R)$ is unimodal with $\max_i h_i-h_{\lceil s/2\rceil}\leq1$. In characteristic zero, the generic Artinian reduction of $R$ has the strong Lefschetz property. If $R$ is level, then $h(R)$ is concave and alternatingly increasing. In particular, it is log-concave and has a peak at position $\lceil s/2\rceil$. If $h_1=h_s$, then $R$ is level and $h(R)=(1,h_1,\ldots,h_1)$.
\end{maintheorem}

\vbox{
Theorems~\ref{thm:A} and~\ref{thm:B} suggest the following extension of Conjecture~\ref{conj:classical}.

\begin{conjecture}\label{conj:unimodality}
Every standard graded nearly Gorenstein domain over a field has a unimodal $h$-vector.
\end{conjecture}
}

We work with Frobenius-split Cohen--Macaulay algebras to allow for applications beyond Ehrhart rings. In Example~\ref{ex:rees-ordinary-cubic}, we show that the generic Artinian reduction of a nearly Gorenstein Rees algebra constructed from a smooth ordinary plane cubic over $\mathbb F_2$ has the strong Lefschetz property. For a triangulated ball whose boundary is an induced subcomplex, we recover the alternating inequalities for its $h$-vector and, in characteristic two, the strong Lefschetz property for the generic Artinian reduction of its Stanley--Reisner ring (Corollary~\ref{cor:ball-induced-boundary}).

For an Ehrhart ring, the monomials of the canonical ideal and of its inverse are determined by the lattice points of the cone over $P$. In particular, the hypotheses on the canonical trace ideal in Theorem~\ref{thm:A} are preserved on passing to characteristic two. Since the $h$-vector also depends only on $P$, the resulting inequalities hold over every field.

For a graded algebra $S=\bigoplus_{i\geq0}S_i$ and an integer $p\geq1$, the \emph{$p$th Veronese subring} $S^{(p)}=\bigoplus_{j\geq0}S_{pj}$ is graded by $[S^{(p)}]_j=S_{pj}$. For Theorem~\ref{thm:B}, we show in Section~\ref{sec:domains} that $K[x,y]^{(h_s)}$ maps onto an Artinian reduction $A$ of $R$ over a field extension $K/\kk$. In fact, $A\cong C^{(h_s)}$ for some graded Artinian algebra $C=K[x,y]/J$. By Lemma~\ref{lem:binary-veronese}, the Hilbert function of $C$ is unimodal, so its subsequence $h_i=\dim_K C_{ih_s}$ is unimodal as well. In the level case we also obtain concavity. To estimate the middle entry, we compare the Hilbert series of $R$ and $R/I$ for some canonical ideal $I$. In characteristic zero the Veronese presentation also gives the strong Lefschetz property for the generic Artinian reduction of $R$ (Theorem~\ref{thm:height-unimodality}).

\section{Canonical trace ideals and Artinian reductions}\label{sec:canonical}

In this section we describe the degree-one part of the canonical trace ideal and the multiplication pairing on an Artinian reduction. Let $R$ be a $d$-dimensional standard graded Cohen--Macaulay algebra over $\kk$, and set $s=\deg h_R(t)$. We use the grading convention $[M(a)]_j=M_{j+a}$ and write $\indeg M=\min\{j:M_j\ne0\}$ for the \emph{initial degree} of a nonzero finitely generated graded module $M$.

Let $Q^*(R)$ denote the \emph{graded total ring of fractions} of $R$, obtained by inverting all homogeneous nonzerodivisors. A \emph{graded fractional ideal} is a finitely generated graded $R$-submodule $I\subseteq Q^*(R)$. We write $I^{-1}=(R:I)=\{v\in Q^*(R):vI\subseteq R\}$. If $I$ contains a homogeneous unit of $Q^*(R)$, localization identifies homogeneous maps $I\to R$ with multiplication by homogeneous elements of $I^{-1}$, so $\tr_R(I)=II^{-1}$ \cite[Lemma~2.6]{M26}. A homogeneous ideal is a \emph{canonical ideal} if it is isomorphic, up to a graded shift, to $\omega_R$. Over a domain, a shift of $\omega_R$ can be identified with a graded fractional ideal \cite[Proposition~3.3.18]{BH98}. More generally, if $R$ is reduced and a graded ideal $I\subseteq R$ represents a shift of $\omega_R$, then $I$ contains a homogeneous nonzerodivisor.

In degree one, the product $II^{-1}$ can receive contributions from several degrees of $I$. The next lemma gives conditions under which the lowest-degree component of $I$ suffices.

\begin{lemma}\label{lem:linear-trace-product}
Let $I$ be a graded fractional ideal isomorphic to a shift of $\omega_R$, and suppose that $I$ contains a homogeneous unit of $Q^*(R)$. Set $a=\indeg I$. If $R$ is a domain, $R$ is level, or $\height([\tr_R(\omega_R)]_1R)\ge2$, then
\begin{equation}\label{eq:linear-trace-product}
 [\tr_R(\omega_R)]_1=I_a[I^{-1}]_{1-a}.
\end{equation}
Here the product on the right is the $\kk$-span of the pairwise products.
\end{lemma}
\begin{proof}
We may assume that $\kk$ is infinite. Indeed, both sides of~\eqref{eq:linear-trace-product} commute with extension to $\kk(t)$, since $I$ is finitely presented and $I^{-1}=\Hom_R(I,R)$. Equality descends by faithful flatness.

We have $\tr_R(\omega_R)=II^{-1}$. If $I=uR$, then $I_a=\kk u$ and $I^{-1}=u^{-1}R$, giving~\eqref{eq:linear-trace-product}. This includes $d=0$, when $I=Q^*(R)=R$. We may therefore assume that $d>0$ and $II^{-1}\subseteq\m_R$. If $[II^{-1}]_1=0$, the formula is immediate. Otherwise the trace has initial degree one.

Suppose $R$ is a domain or is level. Since $I$ becomes free of rank one at every minimal prime, $R$ is generically Gorenstein. Hence $I_a$ contains an element with zero annihilator by \cite[Remark~3.14]{M26}, and \cite[Lemma~2.12~(2)]{M26} gives~\eqref{eq:linear-trace-product} in the initial degree of the trace.

It remains to consider $\height([II^{-1}]_1R)\ge2$. Set $X=I_a[I^{-1}]_{1-a}$ and $Y=\sum_{j>a}I_j[I^{-1}]_{1-j}$, so that $[II^{-1}]_1=X+Y$. For $j>a$, the product $I_a[I^{-1}]_{1-j}$ lies in degree $a+1-j\le0$ of the proper ideal $II^{-1}$, so it vanishes. Hence $YI_a=0$.

Since $R$ is Cohen--Macaulay and the field is infinite, we can choose a regular sequence $\theta_1,\theta_2$ in $[II^{-1}]_1$. Since $YI_a=0$, we have $XY=0$. Choose $x_i\in X$ and $y_i\in Y$ with $\theta_i=x_i+y_i$. Then $\theta_2x_1=x_2x_1=\theta_1x_2$. Since $\theta_2$ is regular modulo $\theta_1$, we have $x_1\in\theta_1R$. Both are linear forms, so $x_1=\lambda\theta_1$ for some scalar $\lambda$. If $\lambda=0$, then $\theta_1=y_1$ annihilates $I_a\ne0$. This contradicts the regularity of $\theta_1$. Hence $x_1$ is a nonzerodivisor. Since $x_1Y=0$, we have $Y=0$. Thus $[II^{-1}]_1=X$, as required.
\end{proof}

For a sequence $\Theta=(\theta_1,\ldots,\theta_n)$ and an $R$-module $N$, write $\Theta N=(\theta_1,\ldots,\theta_n)N$. A sequence $\Theta$ of $d$ linear forms is a \emph{linear system of parameters} if $R/\Theta R$ is Artinian. The algebra $R/\Theta R$ is an \emph{Artinian reduction} of $R$. We allow a field extension when choosing the parameters.

Fix a $\kk$-basis $x_1,\ldots,x_e$ of $R_1$. Let $F=\kk(t_{ij}:1\le i\le d,\ 1\le j\le e)$, where the $t_{ij}$ are algebraically independent over $\kk$, and set $\theta_i=\sum_{j=1}^e t_{ij}x_j$ and $R_F=R\otimes_{\kk}F$. Following Papadakis and Petrotou \cite[Definition~2.2]{PP20}, we call $A_{\mathrm{gen}}=R_F/(\theta_1,\ldots,\theta_d)R_F$ the \emph{generic Artinian reduction} of $R$. The forms $\theta_1,\ldots,\theta_d$ are a linear system of parameters, since such systems form a nonempty open subset of the coefficient space over $\overline{\kk}$.

A standard graded Artinian algebra $A$ has the \emph{strong Lefschetz property} if there is a linear form $\ell\in A_1$ such that, for every pair of integers $0\le a<b$, the map $\times\ell^{b-a}\colon A_a\to A_b$ has maximal rank. Such a form is called a \emph{strong Lefschetz element}.

By \cite[Corollary~4.4.6~(a)]{BH98}, $\indeg\omega_R=d-s$ and $\dim_{\kk}[\omega_R]_{d-s}=h_s$.

To compare entries of $h(R)$, we reduce $R$ and $\omega_R$ modulo the same parameters. The Hilbert function of the first reduction gives $h_0,\ldots,h_s$ starting in degree zero, and that of the second gives $h_s,\ldots,h_0$ starting in degree $d-s$.

\begin{lemma}\label{lem:reduction-pairing}
Let $K/\kk$ be any field extension and let $\Theta$ be any linear system of parameters of $R_{K}=R\otimes_{\kk}K$. Set $A=R_{K}/\Theta R_{K}$ and $M=(\omega_R\otimes_{\kk}K)/\Theta (\omega_R\otimes_{\kk}K)$. For each integer $i$, multiplication induces a perfect pairing $A_i\times M_{d-i}\longrightarrow M_d\cong K$, that is, a $K$-linear isomorphism $M_{d-i}\to\Hom_K(A_i,M_d)$ given by $m\mapsto(a\mapsto am)$. The highest nonzero degree of $M$ is $d$, and $\dim_{K}A_i=\dim_{K}M_{d-i}=h_i$ for every $i$ with $0\le i\le s$. Every nonzero class in $A_i$ pairs nontrivially with the image of some element of $[\omega_R]_{d-i}$. Likewise, every nonzero class in $M_{d-i}$ pairs nontrivially with the image of some element of $R_i$.
\end{lemma}
\begin{proof}
The sequence $\Theta$ is regular because $R_{K}$ is Cohen--Macaulay. Consequently $\Hilb_A(t)=h_R(t)$ and $M(d)\cong\omega_A\cong\Hom_{K}(A,K)$ as graded $A$-modules \cite[Corollary~3.6.14, Theorem~3.6.19~(a), and Exercise~3.3.31]{BH98}. The component of $\Hom_K(A,K)$ in degree $-i$ is $\Hom_K(A_i,K)$, and $A$ acts by $(b\varphi)(a)=\varphi(ba)$. Thus $M_{d-i}$ is the dual of $A_i$, and the pairing is evaluation since $(b\varphi)(1)=\varphi(b)$. Since $A_0=K$ and $A_i=0$ for every $i<0$, the highest nonzero degree of $M$ is $d$. Finally, the images of $\kk$-bases of $[\omega_R]_{d-i}$ and $R_i$ span $M_{d-i}$ and $A_i$, respectively, over $K$. For each nonzero class, at least one such basis element therefore has nonzero product with it.
\end{proof}

Under $(\natur)$ we may choose the parameters in $[\tr_R(\omega_R)]_1$ after extending $\kk$ to $\kk(t)$ if necessary \cite[Remark~3.3~(1), (3), (4)]{M26}. The highest nonzero degrees of the two reductions in Lemma~\ref{lem:reduction-pairing} then give the following inclusions.

\begin{lemma}\label{lem:trace-degree-bounds}
If $R$ satisfies $(\natur)$, then $R_{>s}\subseteq[\tr_R(\omega_R)]_1R$ and $[\omega_R]_{>d}\subseteq[\tr_R(\omega_R)]_1\omega_R$.
\end{lemma}
\begin{proof}
Choose such a system of parameters $\Theta$. By Lemma~\ref{lem:reduction-pairing}, $R/\Theta R$ and $\omega_R/\Theta\omega_R$ vanish in degrees greater than $s$ and $d$, respectively. Each $\theta_i$ belongs to $[\tr_R(\omega_R)]_1$, so the inclusions follow. They descend to $\kk$ by faithful flatness.
\end{proof}

The following lemma gives the two inclusions used in the proof of Theorem~\ref{thm:A}.

\begin{lemma}\label{lem:canonical-squares}
Let $I\subsetneq R$ be a graded ideal containing a homogeneous nonzerodivisor and isomorphic to $\omega_R$ as a graded module. Set $r=\indeg I$, so that $s=d-r$. Suppose that at least one of the following conditions holds.
\begin{enumerate}
\item $R$ satisfies $(\natur)$;
\item $R$ is level and anticanonical level, and $[\tr_R(\omega_R)]_1\ne0$;
\item $R$ is level and $r=1$.
\end{enumerate}
Then
\[
\begin{aligned}
 u^2&\in I_rR &&\text{for every }i\text{ with }0\leq 2i\leq s\text{ and every }u\in I_{d-i},\\
 u^2&\in [I^{-1}]_{1-r}R &&\text{for every }i\text{ with }0\leq 2i<s\text{ and every }u\in R_{s-i}.
\end{aligned}
\]
\end{lemma}
\begin{proof}
Since $I$ is proper, $r\ge1$. If $s=0$, then $d=r$ and only the first inclusion occurs, with $u\in I_r$, so $u^2\in I_rR$. We may therefore assume $s\ge1$, in which case $d=s+r\ge2$.

If $R$ satisfies $(\natur)$, then $\sqrt{[\tr_R(I)]_1R}=\m_R$, so $[\tr_R(I)]_1R$ has height $d\geq2$. By Lemmas~\ref{lem:linear-trace-product} and~\ref{lem:trace-degree-bounds}, we have $[\tr_R(I)]_1=I_r[I^{-1}]_{1-r}$, $I_{>d}\subseteq[\tr_R(I)]_1I$, and $R_{>s}\subseteq[\tr_R(I)]_1R$. If $0\leq2i\leq s$, then every $u\in I_{d-i}$ has square in $I_{2(d-i)}$, where $2(d-i)\geq 2d-s=d+r>d$. Hence $u^2\in[\tr_R(I)]_1I\subseteq I_rR$, since $I^{-1}I\subseteq R$. If $0\leq2i<s$, then $2(s-i)\geq s+1$. Thus every $u\in R_{s-i}$ satisfies $u^2\in[\tr_R(I)]_1R\subseteq[I^{-1}]_{1-r}R$, since $I_r\subseteq R$.

In case~\textup{(2)}, a Gorenstein ring satisfies $(\natur)$, so we may assume that $II^{-1}$ is proper. The generating degrees of $I$ and $I^{-1}$ sum to one, since their product is generated in a single degree and $[II^{-1}]_1\ne0$. Hence $I=I_rR$ and $I^{-1}=[I^{-1}]_{1-r}R$. The first inclusion follows because $u^2\in I$ for every $u\in I$. The second follows from $R\subseteq I^{-1}$.

Finally, under condition~\textup{(3)}, we have $r=1$, $I=I_1R$, and $1\in[I^{-1}]_0$. Hence $u^2\in I_1R$ for every $u\in I$, while $R\subseteq[I^{-1}]_0R$.
\end{proof}

\section{Frobenius splittings and Ehrhart rings}\label{sec:frobenius}

In this section, we prove the inequalities of Theorem~\ref{thm:A}: $h_i\leq h_{s-i}$ for every $i$ with $0\leq 2i\leq s$, and $h_{s-i}\leq h_{i+1}$ for every $i$ with $0\leq 2i<s$. The Frobenius splitting of an Ehrhart ring in characteristic two motivates a criterion for these inequalities in terms of two multiplication maps on an Artinian reduction.

In characteristic two, a \emph{Frobenius splitting} of a ring $R$ is an additive map $\sigma:R\to R$ such that $\sigma(a^2b)=a\sigma(b)$ for all $a,b\in R$ and $\sigma(1)=1$. An ideal $I$ is \emph{compatibly split} if $\sigma(I)\subseteq I$ \cite[Definition~1.1.3]{BK05}. If $a^{2^e}\in I$ for some $e\geq0$, then repeated application of $\sigma$ gives $a\in I$. Thus every compatibly split ideal is radical, and taking $I=0$ shows that $R$ is reduced.

\begin{example}\label{ex:ehrhart-splitting}
Let $P\subset\RR^q$ be a lattice polytope and take $\kk=\mathbb F_2$. Set $C=\RR_{\geq0}(P\times\{1\})$ and $L=\operatorname{span}_{\RR}(C)\cap\ZZ^{q+1}$. The Ehrhart ring is $A_P=\mathbb F_2[C\cap L]$. On its monomial basis define
\[
 \sigma_P(X^{\boldsymbol a})=
 \begin{cases}X^{\boldsymbol a/2},&\boldsymbol a\in2L,\\0,&\boldsymbol a\notin2L.\end{cases}
\]
If $\boldsymbol a\in C\cap2L$, then $\boldsymbol a/2\in C\cap L$. Since $\sigma_P(1)=1$ and $\sigma_P(X^{2\boldsymbol b+\boldsymbol a})=X^{\boldsymbol b}\sigma_P(X^{\boldsymbol a})$, the $\mathbb F_2$-linear extension of $\sigma_P$ is a Frobenius splitting. The monomials $X^{\boldsymbol a}$ with $\boldsymbol a\in\relint(C)\cap L$ span a proper ideal $I$. By the Danilov--Stanley theorem \cite[Theorem~6.3.5]{BH98}, $I\cong\omega_{A_P}$ as graded $A_P$-modules for the Ehrhart grading. This ideal is compatibly split, since $\boldsymbol a/2\in\relint(C)\cap L$ whenever $\boldsymbol a\in\relint(C)\cap2L$.
\end{example}

\begin{example}\label{ex:stanley-reisner-splitting}
For a simplicial complex $\Delta$ on the vertices $1,\ldots,n$, set $S=\mathbb F_2[x_1,\ldots,x_n]$. On the monomials of $S$, define
\[
\sigma(x^{\boldsymbol a})=\begin{cases}
x^{\boldsymbol a/2}, & \text{if every exponent is even},\\
0, & \text{if some exponent is odd}.
\end{cases}
\]
Multiplying by the square of a monomial adds even exponents, so the $\mathbb F_2$-linear extension satisfies $\sigma(u^2m)=u\sigma(m)$ for all $u,m\in S$. Since $\sigma(1)=1$, this is a Frobenius splitting of $S$. For every monomial $m$ with $\sigma(m)\ne0$, the same variables divide $m$ and $\sigma(m)$. Thus $\sigma$ preserves every ideal generated by squarefree monomials, including the Stanley--Reisner ideal of $\Delta$. The induced Frobenius splitting of $\mathbb F_2[\Delta]$ preserves every ideal generated by squarefree monomials for the same reason.
\end{example}

The following is the key lemma in the proof of Theorem~\ref{thm:A}. We take the coefficients of the linear combinations to be independent variables so that we can differentiate with respect to them. This gives a relation for the square of a product. Applying the Frobenius splitting then gives a relation for the product itself.

\begin{lemma}\label{lem:frobenius}
Let $R$ be a standard graded algebra over a field $\kk$ of characteristic two. Let $\sigma$ be a Frobenius splitting of $R$, and let $\mathfrak a$ be an ideal with $\sigma(\mathfrak a)\subseteq\mathfrak a$. Let $U\subseteq Q^*(R)$ be a nonzero finite-dimensional $\kk$-subspace consisting of homogeneous elements of degree $\delta$. Choose bases $x_1,\ldots,x_e$ of $R_1$ and $g_1,\ldots,g_b$ of $U$, and fix $n\geq1$. Let $T$ denote a tuple of algebraically independent variables over $\kk$ containing the entries of an $n\times e$ matrix $(t_{ij})$ and the variables $z_1,\ldots,z_b$. Over $K=\kk(T)$, set $\theta_i=\sum_{j=1}^{e}t_{ij}x_j$, $q=\sum_{\ell=1}^{b}z_\ell g_\ell$, and $\Theta=(\theta_1,\ldots,\theta_n)$. Here $R_K=R\otimes_{\kk}K$ and $\mathfrak a_K=\mathfrak a\otimes_{\kk}K$.

Let $f\in R_K$, and let $u$ be any nonzero homogeneous element of $Q^*(R)$. If $u^2\in UR$ and $2\deg u-\delta\geq n$, then $uf\in\Theta \mathfrak a_K$ whenever $qf^2\in\Theta \mathfrak a_K$.
\end{lemma}
\begin{proof}
By hypothesis, $qf^2=\sum_{i=1}^n\theta_i f_i$ for some $f_1,\ldots,f_n\in\mathfrak a_K$. Multiplying by $\theta_1\cdots\theta_n$ gives
\begin{equation}\label{eq:derivative-squares}
 q\theta_1\cdots\theta_n f^2=\sum_{i=1}^n\theta_i^2 F_i,\qquad F_i=\Bigl(\prod_{j\ne i}\theta_j\Bigr)f_i\in\mathfrak a_K.
\end{equation}
For each variable $v$ in $T$, let $\partial_v$ denote formal differentiation with respect to $v$ on $Q^*(R)\otimes_{\kk}K$. It sends $v$ to $1$ and vanishes on $Q^*(R)$ and on the other variables in $T$. The product rule determines $\partial_v$ on rational functions as well, since $\partial_v(h^{-1})=-h^{-2}\partial_v h$ for every nonzero $h\in\kk[T]$. If $c_\nu\in K$ and $a_\nu\in\mathfrak a$, then $\partial_v(\sum_\nu c_\nu a_\nu)=\sum_\nu(\partial_v c_\nu)a_\nu\in\mathfrak a_K$, since $\partial_v c_\nu\in K$. In characteristic two, $\partial_v(f^2)=2f\,\partial_v f=0$, even if $f$ depends on $v$.

Fix $1\leq\ell\leq b$ and $j_1,\ldots,j_n\in\{1,\ldots,e\}$. Apply $D=\partial_{z_\ell}\partial_{t_{1j_1}}\cdots\partial_{t_{nj_n}}$ to~\eqref{eq:derivative-squares}. On the left, the derivatives annihilate $f^2$. Among the remaining factors, $z_\ell$ occurs only in $q$, with $\partial_{z_\ell}q=g_\ell$, and $t_{ij_i}$ occurs only in $\theta_i$, with $\partial_{t_{ij_i}}\theta_i=x_{j_i}$.

On the right, each coefficient derivative $\partial$ satisfies $\partial(\theta_i^2)=2\theta_i\partial(\theta_i)=0$ in characteristic two. The product rule therefore gives $\partial(\theta_i^2 F)=\theta_i^2\partial F$ for every $F$. Applying this at each step yields $D(\theta_i^2 F_i)=\theta_i^2D(F_i)$, so~\eqref{eq:derivative-squares} becomes
\begin{equation}\label{eq:derivative-products}
 g_\ell x_{j_1}\cdots x_{j_n}f^2=\sum_{i=1}^n\theta_i^2D(F_i)\in(\theta_1^2,\ldots,\theta_n^2)\mathfrak a_K.
\end{equation}

By homogeneity, $u^2\in UR$ gives $u^2=\sum_{\ell=1}^b g_\ell w_\ell$ with $w_\ell\in R_{2\deg u-\delta}$. Since $R$ is standard graded and $2\deg u-\delta\geq n$, each $w_\ell$ is an $R$-linear combination of the products $x_{j_1}\cdots x_{j_n}$. Taking the corresponding linear combinations in~\eqref{eq:derivative-products} gives
\[
 (uf)^2=\sum_{\ell=1}^b g_\ell w_\ell f^2\in(\theta_1^2,\ldots,\theta_n^2)\mathfrak a_K.
\]

We now use the splitting to pass from $(uf)^2$ to $uf$. On $R[T]$, set
\[
 \widetilde\sigma\!\left(\sum_{\nu}a_{\nu}T^{\nu}\right)=\sum_{\mu}\sigma(a_{2\mu})T^\mu.
\]
Thus we keep the terms whose exponents are all even and halve their exponents. Multiplication by $c^2T^{2\alpha}$ gives the factor $cT^\alpha$ after applying $\widetilde\sigma$. Since squares add in characteristic two and $\widetilde\sigma(1)=1$, this defines a splitting. It preserves $\mathfrak a[T]$ because $\sigma(\mathfrak a)\subseteq\mathfrak a$. The localization formula $\widetilde\sigma(h/s)=\widetilde\sigma(hs)/s$ extends the splitting to $R_K$. This extension preserves $\mathfrak a_K$. Localizing further at the homogeneous nonzerodivisors of $R$ gives a splitting of $Q^*(R)\otimes_{\kk}K$ that agrees with it on $R_K$. Choose $b_i\in\mathfrak a_K$ such that $(uf)^2=\sum_i\theta_i^2 b_i$. Applying $\widetilde\sigma$ gives
\[
 uf=\widetilde\sigma\bigl((uf)^2\bigr)=\sum_{i=1}^n\theta_i\widetilde\sigma(b_i)\in\Theta \mathfrak a_K.\qedhere
\]
\end{proof}

In characteristic two, Adiprasito, Papadakis, and Petrotou differentiate the volume map with respect to the coefficients of a linear system of parameters \cite[Lemma~5.1]{APP25}. They deduce square nonvanishing and pass to characteristic zero by reduction modulo two \cite[Theorems~3.3 and~6.5]{APP25}, as in \cite[Section~3.2]{KLS26}. In Lemma~\ref{lem:frobenius} we differentiate an ideal membership relation, so the elements $g_\ell$ in the expression for $q$ may belong to a fractional ideal.

Let $I$ be a graded ideal containing a nonzerodivisor. After reduction modulo the parameters, multiplication by an element of $I$ gives a map from the reduction of $R$ to that of $I$. Multiplication by an element of $I^{-1}$ gives a map in the opposite direction. We use Lemma~\ref{lem:frobenius} to prove injectivity for the two maps below.

\begin{proposition}\label{prop:split-injections}
Let $R$ be a standard graded algebra over a field $\kk$ of characteristic two, with a Frobenius splitting $\sigma$. Let $I$ be a graded ideal containing a homogeneous nonzerodivisor, and suppose that $\sigma(I)\subseteq I$. Set $d=\dim R$, $r=\indeg I$, and $s=d-r$. Choose $\kk$-bases $x_1,\ldots,x_e$ of $R_1$, $g_1,\ldots,g_b$ of $I_r$, and $v_1,\ldots,v_c$ of $[I^{-1}]_{1-r}$. Let $F=\kk(t_{j\ell})$ and $K=F(z_1,\ldots,z_b,w_1,\ldots,w_c)$, where all the coefficients are algebraically independent over $\kk$. For each $j$ with $1\leq j\leq d$, set $\theta_j=\sum_{\ell=1}^e t_{j\ell}x_\ell$. Also set
\[
 q=\sum_{\ell=1}^b z_\ell g_\ell,\qquad q'=\sum_{\ell=1}^c w_\ell v_\ell.
\]
If $[I^{-1}]_{1-r}=0$, take $q'=0$. Set $R_K=R\otimes_{\kk}K$ and $I_K=I\otimes_{\kk}K$. The forms $\Theta=(\theta_1,\ldots,\theta_d)$ are a linear system of parameters of $R_K$. Set $A=R_K/\Theta R_K\cong A_{\mathrm{gen}}\otimes_F K$ and $M=I_K/\Theta I_K$. In each assertion $f$ is any homogeneous representative of the class $\bar f$ over $K$, and $u$ is chosen over $\kk$.
\begin{enumerate}
\item Let $0\leq2i\leq s$. Suppose that for every $0\neq\bar f\in A_i$ there is an element $u\in I_{d-i}$ such that $u\bar f\neq0$ in $M_d$ and $u^2\in I_r R$. Then $qf^2\notin\Theta I_K$ for every such $\bar f$. In particular, multiplication by $q$ gives an injection $A_i\longrightarrow M_{r+i}$.
\item Let $0\leq2i<s$. Suppose that for every $0\neq\bar f\in M_{r+i}$ there is an element $u\in R_{s-i}$ such that $u\bar f\neq0$ in $M_d$ and $u^2\in[I^{-1}]_{1-r}R$. Then $q'f^2\notin\Theta I_K$ for every such $\bar f$. In particular, multiplication by $q'$ gives an injection $M_{r+i}\longrightarrow A_{i+1}$.
\end{enumerate}
In particular, if both hypotheses hold at $i$, then $\dim_K A_i\leq\dim_K M_{r+i}\leq\dim_K A_{i+1}$.
\end{proposition}
\begin{proof}
If $d=0$, then $R=I=\kk$, and the assertions follow from $qf^2\ne0$ for nonzero scalars $q$ and $f$. Suppose $d\geq1$. The $d$-tuples of linear parameters form a nonempty open set over $\overline{\kk}$. Since the $t_{j\ell}$ are algebraically independent, $\Theta$ is a system of parameters. The two multiplication maps are well defined because $qR_K\subseteq I_K$ and $q'I_K\subseteq R_K$.

For \textup{(1)}, let $0\neq\bar f\in A_i$ and choose a representative $f\in[R_K]_i$. Suppose $qf^2\in\Theta I_K$ and choose $u$ as in \textup{(1)}. Since $2i\leq s$, we have $2\deg u-r=2(d-i)-r=d+s-2i\geq d$. Together with $u^2\in I_rR$, this allows us to apply Lemma~\ref{lem:frobenius} with $\mathfrak a=I$ and $U=I_r$. It follows that $uf\in\Theta I_K$. Thus $u\bar f=0$ in $M_d$, contrary to the choice of $u$.

For \textup{(2)}, let $0\neq\bar f\in M_{r+i}$ and choose a representative $f\in[I_K]_{r+i}$. Suppose $q'f^2\in\Theta I_K$ and choose $u$ as in \textup{(2)}. Its nonzero square belongs to $[I^{-1}]_{1-r}R$, so $[I^{-1}]_{1-r}\neq0$. Since $2i\leq s-1$, we have $2\deg u-(1-r)=2(s-i)-1+r=d+s-1-2i\geq d$. Apply Lemma~\ref{lem:frobenius} with $\mathfrak a=I$, $U=[I^{-1}]_{1-r}$, and $q'$ in place of $q$. We again get $u f\in\Theta I_K$, contrary to the choice of $u$.

Finally, if $q\bar f=0$ in~\textup{(1)}, multiply $qf\in\Theta I_K$ by $f\in R_K$. If $q'\bar f=0$ in~\textup{(2)}, multiply $q'f\in\Theta R_K$ by $f\in I_K$. Both contradict the respective square nonvanishing.

More generally, the proof of \textup{(1)} applies in any degree if, for each nonzero source class, one can choose $u$ with nonzero product in $M$, $u^2\in I_rR$, and $2\deg u-r\ge d$. The proof of \textup{(2)} applies under the analogous assumptions $u^2\in[I^{-1}]_{1-r}R$ and $2\deg u-(1-r)\ge d$.
\end{proof}

For a Cohen--Macaulay algebra, the pairing in Lemma~\ref{lem:reduction-pairing} gives the nonzero products required in Proposition~\ref{prop:split-injections}. We obtain the following criterion.

\begin{theorem}\label{thm:split-canonical}
Let $R$ be a standard graded Cohen--Macaulay algebra of dimension $d$ over a field $\kk$ of characteristic two. Let $I\subseteq R$ be a graded ideal isomorphic to $\omega_R$ as a graded module, and suppose that $R$ admits a Frobenius splitting $\sigma$ with $\sigma(I)\subseteq I$. Let $r=\indeg I$ and $h(R)=(h_0,\ldots,h_s)$, so that $s=d-r$.

For a fixed integer $i$, consider the conditions
\begin{alignat}{2}
 u^2&\in I_rR &&\qquad\text{for every }u\in I_{d-i},\label{eq:canonical-square-membership}\\
 u^2&\in [I^{-1}]_{1-r}R &&\qquad\text{for every }u\in R_{s-i}.\label{eq:inverse-square-membership}
\end{alignat}
If $0\leq 2i\leq s$ and \eqref{eq:canonical-square-membership} holds, then $h_i\leq h_{s-i}$. If $0\leq 2i<s$ and \eqref{eq:inverse-square-membership} holds, then $h_{s-i}\leq h_{i+1}$. If both conditions hold for every $i$ in their respective ranges, then $h(R)$ is alternatingly increasing, and hence unimodal.
\end{theorem}
\begin{proof}
Since $R$ is reduced, its canonical ideal $I$ contains a homogeneous nonzerodivisor. Take $A$ and $M$ as in Proposition~\ref{prop:split-injections}. Lemma~\ref{lem:reduction-pairing} gives, for every nonzero class in $A_i$, an element $u\in I_{d-i}$ pairing nontrivially with it. Likewise, each nonzero class in $M_{r+i}$ pairs nontrivially with some $u\in R_{s-i}$. Conditions~\eqref{eq:canonical-square-membership} and~\eqref{eq:inverse-square-membership}, respectively, now give the two injections of Proposition~\ref{prop:split-injections}. Their dimensions give the desired inequalities, since $\dim_K A_j=h_j$ and $\dim_K M_{r+i}=h_{s-i}$.
\end{proof}

\begin{proof}[Proof of Theorem~\ref{thm:A}]
For the cone $C$ over $P\times\{1\}$, let $I$ denote the ideal of interior monomials in $A_P$. The exponent sets of $I$ and $I^{-1}$ depend only on the lattice points of $C$ \cite[Proposition~4.1 and Corollary~4.2]{HMP19}. Thus the degrees of their minimal monomial generators, and hence levelness and anticanonical levelness, are independent of the field. We have $[II^{-1}]_1\neq0$ precisely when $II^{-1}$ contains a degree-one monomial. Moreover, $(\natur)$ holds precisely when for every degree-one monomial $m$ of $A_P$ there is some $N\geq1$ such that $m^N\in[II^{-1}]_1A_P$. Both statements depend only on the exponent sets of $I$ and $I^{-1}$. The Hilbert series counts lattice points, so we may work over $\mathbb F_2$. By Example~\ref{ex:ehrhart-splitting}, the monomial splitting $\sigma_P$ preserves $I$. In case~\textup{(3)}, an interior lattice point of $P$ gives $\indeg I=1$. Thus in each case Lemma~\ref{lem:canonical-squares} gives \eqref{eq:canonical-square-membership} and \eqref{eq:inverse-square-membership} in their respective ranges. The alternating inequalities now follow from Theorem~\ref{thm:split-canonical}.
\end{proof}

\begin{remark}
By \cite[Proposition~4.1]{M26}, the Ehrhart ring of an IDP lattice polytope $P$ satisfies $(\natur)$ exactly when its canonical trace ideal contains $\mathbf{x}^w t$ for every vertex $w$ of $P$.

We call $R$ \emph{locally Gorenstein} if $R_{\mathfrak p}$ is Gorenstein for every prime $\mathfrak p\ne\m_R$. Every ring satisfying $(\natur)$ is locally Gorenstein because the canonical trace ideal defines the non-Gorenstein locus \cite[Lemma~2.1]{HHS19}. Even a locally Gorenstein Ehrhart ring can have a non-unimodal $h$-vector. Ferroni constructed a smooth IDP polytope $P$ with non-unimodal $h^*(P)$ \cite[Theorem~4.1]{F26}. Since $A_P$ is standard graded and $\operatorname{Proj} A_P$ is smooth, $A_P$ is regular on the punctured spectrum.
\end{remark}

Under the hypotheses for the alternating inequalities in Theorem~\ref{thm:split-canonical}, one can choose $q$ and $q'$ so that $\overline{qq'}$ is a strong Lefschetz element of the generic Artinian reduction.

\begin{theorem}\label{thm:split-strong-lefschetz}
Let $R$ be a $d$-dimensional standard graded Cohen--Macaulay algebra over a field $\kk$ of characteristic two. Let $I\subseteq R$ be a graded ideal isomorphic to $\omega_R$ as a graded $R$-module. Suppose that $R$ has a Frobenius splitting $\sigma$ such that $\sigma(I)\subseteq I$. Set $r=\indeg I$ and $s=d-r$. Assume that
\[
\begin{aligned}
u^2&\in I_rR &&\text{for every $i$ with $0\leq 2i\leq s$ and every $u\in I_{d-i}$},\\
u^2&\in[I^{-1}]_{1-r}R &&\text{for every $i$ with $0\leq 2i<s$ and every $u\in R_{s-i}$}.
\end{aligned}
\]
Let $A_{\mathrm{gen}}$ be the generic Artinian reduction of $R$ over $F=\kk(t_{ij})$, where the $t_{ij}$ are the coefficients of the generic linear system of parameters. In fixed $\kk$-bases of $I_r$ and $[I^{-1}]_{1-r}$, respectively, take linear combinations $q$ and $q'$ whose coefficients are jointly algebraically independent over $F$. Let $K$ be the field obtained by adjoining these coefficients to $F$. Then the image of $qq'$ is a strong Lefschetz element of $A=A_{\mathrm{gen}}\otimes_F K$. Moreover, one can choose the coefficients of $q$ and $q'$ in $F$ so that the image of their product is a strong Lefschetz element of $A_{\mathrm{gen}}$.
\end{theorem}
\begin{proof}
If $s=0$, then $h_R(t)=1$, so both $A_{\mathrm{gen}}$ and $A$ are concentrated in degree zero. For every choice of coefficients, the degree-one element $\overline{qq'}$ is therefore zero. It is a strong Lefschetz element since every map in the definition has zero target and hence maximal rank. Now assume $s\ge1$. Let $\Theta$ be the generic linear system of parameters defining $A_{\mathrm{gen}}$, and set $I_K=I\otimes_{\kk}K$ and $M=I_K/\Theta I_K$. Consider the sequence of multiplication maps
\[
 A_0\xrightarrow{\times q}M_r\xrightarrow{\times q'}A_1\longrightarrow\cdots\longrightarrow A_s\xrightarrow{\times q}M_d.
\]
Set $D_{2i}=A_i$ and $D_{2i+1}=M_{r+i}$ for each $i$ with $0\le i\le s$. Alternating multiplication by $q$ and $q'$ defines a degree-one linear map $\varphi$ on $\bigoplus_{j=0}^{2s+1}D_j$, with $\varphi(D_{2s+1})=0$. For every $0\le j\le 2s+1$, write $\langle x,y\rangle=xy$ for the perfect pairing $D_j\times D_{2s+1-j}\to M_d$ of Lemma~\ref{lem:reduction-pairing}. For every $0\le j\le 2s$ and all $x\in D_j$, $y\in D_{2s-j}$, we have $\langle\varphi x,y\rangle=\langle x,\varphi y\rangle$. If $j$ is even, $x,y\in A$ and both sides equal $\bar q(xy)$, where $\bar q$ is the class of $q$ in $M_r$. If $j$ is odd, take representatives $\tilde x,\tilde y\in I_K$ of $x,y\in M$. The two sides are $(q'x)y$ and $x(q'y)$, respectively. Since $q'I_K\subseteq R_K$, both are the class of $(q'\tilde x)\tilde y=\tilde x(q'\tilde y)$ modulo $\Theta I_K$.

For $0\le 2i\le s$, Lemma~\ref{lem:reduction-pairing} gives, for every nonzero $x\in A_i$, an element $u\in I_{d-i}$ such that $\bar u x\ne0$ in $M_d$. For $0\le 2i<s$, the same lemma gives, for every nonzero $y\in M_{r+i}$, an element $v\in R_{s-i}$ such that $\bar v y\ne0$ in $M_d$, since $d-(r+i)=s-i$. Here $\bar u$ and $\bar v$ are the images in $M$ and $A$, respectively. The hypotheses on squares give $u^2\in I_rR$ and $v^2\in[I^{-1}]_{1-r}R$ in the respective ranges. Proposition~\ref{prop:split-injections}\,\textup{(1)} and \textup{(2)} therefore give the injectivity of $\varphi\colon D_j\to D_{j+1}$ for every $0\le j\le s$.

The middle map $D_s\to D_{s+1}$ is $A_i\xrightarrow{\times q}M_{r+i}$ when $s=2i$, and $M_{r+i}\xrightarrow{\times q'}A_{i+1}$ when $s=2i+1$. Take any nonzero $z\in D_s$ and choose a homogeneous representative $f$. The element $\langle z,\varphi z\rangle\in M_d$ is represented by $qf^2$ if $s$ is even and by $q'f^2$ if $s$ is odd. Proposition~\ref{prop:split-injections}\,\textup{(1)} gives $qf^2\notin\Theta I_K$ when $s$ is even, and Proposition~\ref{prop:split-injections}\,\textup{(2)} gives $q' f^2\notin\Theta I_K$ when $s$ is odd. Thus $\langle z,\varphi z\rangle\ne0$.

Fix $j$ with $0\le j\le s$. For every nonzero $x\in D_j$, we have $z=\varphi^{s-j}x\ne0$ and
\[
 \langle x,\varphi^{2s+1-2j}x\rangle=\langle\varphi^{s-j}x,\varphi^{s+1-j}x\rangle=\langle z,\varphi z\rangle\ne0.
\]
Thus $\varphi^{2s+1-2j}:D_j\to D_{2s+1-j}$ is injective, and hence an isomorphism by duality. Now fix any pair $0\le m\le n\le 2s+1$. If $m+n\le2s+1$, the isomorphism starting in $D_m$ factors through $\varphi^{n-m}\colon D_m\to D_n$, so this map is injective. If $m+n>2s+1$, the isomorphism from $D_{2s+1-n}$ to $D_n$ factors through $\varphi^{n-m}\colon D_m\to D_n$, so this map is surjective.

On $A_a=D_{2a}$, the map $\varphi^{2(b-a)}$ is multiplication by $\ell^{b-a}$, where $\ell=\overline{qq'}$. The map $A_a\to A_b$ is injective if $a+b\le s$ and surjective if $a+b\ge s+1$. Thus it has rank $\min\{h_a,h_b\}$ for every $0\le a<b\le s$, and $\ell$ is a strong Lefschetz element over $K$. In fixed $F$-bases, choose one nonzero maximal minor for each of these maps. Their product is a nonzero polynomial over $F$ in the coefficients of $q$ and $q'$. Since $F$ is infinite, we may choose the coefficients in $F$ so that this product is nonzero. The resulting $\overline{qq'}$ is a strong Lefschetz element of $A_{\mathrm{gen}}$.
\end{proof}

Combining Lemma~\ref{lem:canonical-squares} with Theorems~\ref{thm:split-canonical} and~\ref{thm:split-strong-lefschetz} gives the following corollary.

\begin{corollary}\label{cor:split-applications}
Let $R$ be a standard graded Cohen--Macaulay algebra over a field of characteristic two, with a Frobenius splitting $\sigma$ and a proper graded ideal $I\cong\omega_R$ such that $\sigma(I)\subseteq I$, where the isomorphism preserves degrees. If one of the following conditions holds, then $h(R)$ is alternatingly increasing and the generic Artinian reduction of $R$ has the strong Lefschetz property:
\begin{enumerate}
\item $R$ satisfies $(\natur)$;
\item $R$ is level and anticanonical level, and $[\tr_R(\omega_R)]_1\neq0$;
\item $R$ is level and $\indeg\omega_R=1$.
\end{enumerate}
\end{corollary}
\begin{proof}
Since $R$ is reduced, $I$ contains a homogeneous nonzerodivisor. In each of the three cases, Lemma~\ref{lem:canonical-squares} gives \eqref{eq:canonical-square-membership} and \eqref{eq:inverse-square-membership} in their respective ranges. Theorem~\ref{thm:split-canonical} gives the alternating inequalities, and Theorem~\ref{thm:split-strong-lefschetz} gives the strong Lefschetz property.
\end{proof}

We apply the multiplication maps to the Ehrhart ring of an IDP lattice polytope. Recall that an \emph{$O$-sequence} is the Hilbert function of a standard graded algebra over a field.

\begin{corollary}\label{cor:ehrhart-consequences}
Let $P$ be an IDP lattice polytope and let $R=A_P$ be its Ehrhart ring over a field $\kk$. Write $h^*(P)=(h_0,\ldots,h_s)$ and set $h_j=0$ for $j<0$ or $j>s$. Suppose that at least one of the following conditions holds.
\begin{enumerate}
\item $R$ satisfies $(\natur)$;
\item $R$ is level and anticanonical level, and $[\tr_R(\omega_R)]_1\ne0$;
\item $R$ is level and $P$ has an interior lattice point.
\end{enumerate}
Then the following assertions hold.
\begin{enumerate}
\renewcommand{\theenumi}{\alph{enumi}}
\renewcommand{\labelenumi}{\textup{(\theenumi)}}
\item\label{item:ehrhart-lefschetz} If $\kk$ has characteristic zero or two, the generic Artinian reduction of $R$ has the strong Lefschetz property.
\item\label{item:ehrhart-osequences} Both of the following sequences are $O$-sequences.
\[
 \bigl(\max\{h_j-h_{j-1},0\}\bigr)_{j\ge0},\qquad
 \bigl(\max\{h_j-h_{s+1-j},0\}\bigr)_{j\ge0}.
\]
\end{enumerate}
\end{corollary}
\begin{proof}
If $s=0$, the generic Artinian reduction is a field, so $\ell=0$ is a strong Lefschetz element and both sequences are $(1,0,\ldots)$. Assume $s\ge1$.

\smallskip
\noindent\textup{(\ref{item:ehrhart-lefschetz})} Let $I$ be the canonical ideal spanned by the interior monomials as in Example~\ref{ex:ehrhart-splitting}, and set $r=\indeg I$. Use independent variables $\mathbf t$ as coefficients of the parameters $\Theta$ in a monomial basis of $R_1$. Set $F=\kk(\mathbf t)$ and let $A=R_F/\Theta R_F$ be the generic Artinian reduction. In monomial bases of $I_r$ and $[I^{-1}]_{1-r}$, respectively, take linear combinations $q$ and $q'$ with independent coefficients $\mathbf z$ and $\mathbf w$. Set $K=F(\mathbf z,\mathbf w)$ and $A_K=A\otimes_F K$. For $\kk=\mathbb F_2$, Lemma~\ref{lem:canonical-squares} and Theorem~\ref{thm:split-strong-lefschetz} give the strong Lefschetz element $\ell=\overline{qq'}$ of $A_K$.

For each pair $0\le a<b\le s$, let $B_b$ have as columns the products $\theta_j m$ with $m$ a monomial of degree $b-1$, written in the monomial basis of $R_b$. Append the columns of multiplication by $(qq')^{b-a}$ from $R_a$ to $R_b$. The increase in rank is the rank of multiplication by $\overline{qq'}^{\,b-a}$ from $A_a\otimes_F K$ to $A_b\otimes_F K$. In the monomial bases, multiplication adds exponents. Thus the same polynomials in $\mathbf t,\mathbf z,\mathbf w$ with integer coefficients give the entries in every characteristic. Also, $\operatorname{rank}B_b=\dim_{\kk}R_b-h_b$ in every characteristic. For each $(a,b)$ the augmented matrix therefore has a minor of size $\operatorname{rank}B_b+\min\{h_a,h_b\}$ whose determinant is nonzero modulo two. This determinant is nonzero over every field of characteristic zero or two. Since $F=\kk(\mathbf t)$ is infinite, we can choose values in $F$ for $\mathbf z$ and $\mathbf w$ at which the product of the chosen determinants is nonzero. With these coefficients, $\overline{qq'}$ is a strong Lefschetz element of $A$ over $F$.

\smallskip
\noindent\textup{(\ref{item:ehrhart-osequences})} As in the proof of Theorem~\ref{thm:A}, the hypotheses and $h^*(P)$ are independent of the ground field. Take $\kk=\mathbb F_2$ and retain $A_K$, $q$, $q'$ and $\ell=\overline{qq'}$ from part~\textup{(\ref{item:ehrhart-lefschetz})}, with the coefficients of $q$ and $q'$ still algebraically independent over $F$. Since $\ell$ is a strong Lefschetz element, $A_K/(\ell)$ has the first Hilbert function.

Set $I_K=I\otimes_{\kk}K$ and $M=I_K/\Theta I_K$. The degree-$j$ component of the homogeneous ideal $q'M$ is the image of multiplication by $q'$ from $M_{r+j-1}$ to $[A_K]_j$. By Lemma~\ref{lem:reduction-pairing}, the source has dimension $h_{s+1-j}$. For $1\le j\le (s+1)/2$, Lemmas~\ref{lem:reduction-pairing} and~\ref{lem:canonical-squares} give the hypotheses of Proposition~\ref{prop:split-injections}\,\textup{(2)} at $i=j-1$, so this map is injective. For $(s+1)/2<j\le s$, Theorem~\ref{thm:A} gives $h_{j-1}\ge h_j$, so multiplication by $\ell$ from $[A_K]_{j-1}$ to $[A_K]_j$ is surjective. Since $q[A_K]_{j-1}\subseteq M_{r+j-1}$, we obtain
\[
 [A_K]_j=\ell[A_K]_{j-1}\subseteq q'M_{r+j-1}\subseteq[A_K]_j.
\]
Thus multiplication by $q'$ has maximal rank in every degree $1\le j\le s$. Since $M_{r-1}=0$, the degree-zero component of $A_K/q'M$ has dimension one. Hence $A_K/q'M$ has the second Hilbert function. Moreover, $\ell\in q'M$, so there is a natural surjection $A_K/(\ell)\to A_K/q'M$.
\end{proof}

\begin{remark}\label{rem:app-scope}
In case~\textup{(3)} of Theorem~\ref{thm:A}, the alternating inequalities asked for in \cite[Question~2.21]{SVL13} follow from \cite[Corollary~2.2 and Theorem~2.4]{APP25}, since $s=\dim P$. The strong Lefschetz property in Corollary~\ref{cor:ehrhart-consequences}\,\textup{(\ref{item:ehrhart-lefschetz})} follows in this case from \cite[Theorems~2.1 and~6.1]{APP25} with $j=1$ in Theorem~6.1.
\end{remark}

\begin{example}\label{ex:second-difference}
Consider the projective monomial curve with coordinate ring
\[
R=\QQ[u^{51},u^{50}v,u^{49}v^2,u^{48}v^3,u^{37}v^{14},u^{36}v^{15},v^{51}],
\]
with all seven generators in degree one. One can check with \texttt{Macaulay2}~\cite{M2} that $R$ is Cohen--Macaulay and $h(R)=(1,5,9,13,12,9,2)$. This vector is concave and alternatingly increasing, with $h_s=2$ and $h_1-h_s=3$. The first sequence in Corollary~\ref{cor:ehrhart-consequences}\,\textup{(\ref{item:ehrhart-osequences})} is the $O$-sequence $(1,4,4,4,0,\ldots)$. The second begins $(1,3,0,1)$, so it cannot be an $O$-sequence. Thus $h(R)$ cannot be the $h^*$-vector of an IDP polytope satisfying the hypotheses of Theorem~\ref{thm:A}. In particular, $R$ is not nearly Gorenstein.
\end{example}

We conclude this section with two applications beyond Ehrhart rings. First we apply Corollary~\ref{cor:split-applications}\,\textup{(1)} to a Rees algebra.

\begin{example}\label{ex:rees-ordinary-cubic}
Let $C$ be the homogeneous coordinate ring of a smooth ordinary plane cubic over $\kk=\mathbb F_2$. For example, homogenizing the Weierstrass equation $y^2+xy=x^3+x$ in \cite[Section~2, (4)]{Sch22} gives
\[
 C=\kk[x,y,z]/(y^2z+xyz+x^3+xz^2).
\]
Set $A=C[x_1,x_2,x_3]$ and $R=A[\m_A T]$, graded by $\deg(aT^j)=\deg_A a$ for homogeneous $a\in\m_A^j$. Since $A$ is a cubic hypersurface in six variables, $R_{\m_R}$ is Cohen--Macaulay and its canonical trace ideal equals its maximal ideal \cite[Example~2.11]{EY26}. The canonical trace ideal is homogeneous and commutes with localization, so $\tr_R(\omega_R)=\m_R$. Thus $R$ is nearly Gorenstein.

Ordinarity gives a homogeneous Frobenius splitting of $C$ \cite[Example~2.3.1]{KKP21}. Write $R$ as the Segre product $A\mathbin{\#}B$, where $B=\kk[u,v]$ \cite[proof of Proposition~6.1]{Si00}. On each of the polynomial rings $\kk[x_1,x_2,x_3]$ and $B$, we use the Frobenius splitting defined in Example~\ref{ex:stanley-reisner-splitting}. The tensor product of these maps with the splitting of $C$ restricts to a splitting of $R$ that preserves $I=(x_1x_2x_3 A)\mathbin{\#}(uvB)$. Since $\omega_C\cong C$, this proper ideal is a graded canonical ideal by \cite[Theorem~4.3.1]{GW78}. Corollary~\ref{cor:split-applications}\,\textup{(1)} now gives the strong Lefschetz property for the generic Artinian reduction of $R$.
\end{example}

We next apply Corollary~\ref{cor:split-applications}\,\textup{(3)} to the Stanley--Reisner rings of triangulated balls with induced boundary.
We write $h(\Delta)=(h_0(\Delta),\ldots,h_{s+1}(\Delta))$ for the \emph{$h$-vector} of an $s$-dimensional simplicial complex $\Delta$, where $h_i(\Delta)$ is the coefficient of $t^i$ in the $h$-polynomial of its Stanley--Reisner ring $\kk[\Delta]$.
The alternating inequalities in the next corollary are known \cite[Section~5.1, (17) and Theorem~5.1]{Ath26}. For the Lefschetz assertion, the injections follow from the argument of Adiprasito and Yashfe~\cite{AY21} explained in \cite[Section~5.1]{Ath26}. The surjections follow by applying the sphere Lefschetz theorem \cite[Corollary~9.10]{PP20} after coning off $\partial\Delta$.

\begin{corollary}\label{cor:ball-induced-boundary}
Let $\Delta$ be a triangulation of an $s$-dimensional ball such that $\partial\Delta$ is an induced subcomplex of $\Delta$. Then $h_{s+1}(\Delta)=0$ and the sequence $h_0(\Delta),\ldots,h_s(\Delta)$ is alternatingly increasing. Moreover, for every field $\kk$ of characteristic two, the generic Artinian reduction of $\kk[\Delta]$ has the strong Lefschetz property.
\end{corollary}
\begin{proof}
For a single vertex, $h(\Delta)=(1,0)$ and the Artinian reduction is a field. For $s\ge1$, the ring $R=\mathbb F_2[\Delta]$ is Cohen--Macaulay. By Hochster's theorem \cite[Theorem~II.7.3]{S96}, the monomials $x_F=\prod_{v\in F}x_v$ for the interior faces $F$ of $\Delta$ generate a graded canonical ideal $I$ of $R$. Since $\partial\Delta$ is induced, every face supported on boundary vertices lies in $\partial\Delta$. Thus each interior face contains an interior vertex, and the variables of the interior vertices generate $I$. Hence $R$ is level with $\indeg\omega_R=1$, and $\deg h_R=\dim R-1=s$.

Since $I$ is generated by variables, the Frobenius splitting of $R$ in Example~\ref{ex:stanley-reisner-splitting} preserves $I$. Corollary~\ref{cor:split-applications}\,\textup{(3)} gives both assertions over $\mathbb F_2$. Extending the parameter field from $\mathbb F_2(\mathbf t)$ to $\kk(\mathbf t)$ gives the Lefschetz assertion over every field $\kk$ of characteristic two.
\end{proof}

\section{A numerical criterion for unimodality}\label{sec:domains}

In this section we prove Theorem~\ref{thm:B}. Over a suitable field extension $K/\kk$, we first construct a graded surjection from a Veronese subring of $K[x,y]$ onto an Artinian reduction of $R$. Its Hilbert function gives unimodality. To estimate the middle entry of $h(R)$, we then use a canonical ideal of $R$.

By Lemma~\ref{lem:linear-trace-product}, the degree-one part of the canonical trace ideal of a standard graded Cohen--Macaulay domain is spanned by products $vu$, where $u$ belongs to the lowest-degree component of a canonical fractional ideal and $v$ to the complementary degree of its inverse. The following lemma uses these products to construct the required surjection.

\begin{lemma}\label{lem:product-matrix}
Let $R$ be a standard graded Cohen--Macaulay domain over a field $\kk$. Let $E,W\subseteq Q^*(R)$ be nonzero finite-dimensional $\kk$-subspaces, each concentrated in a single degree, such that $EW\subseteq R_1$. Suppose that $\dim_{\kk}E\geq2$ and $\height((EW)R)\geq2$. Set $p=\dim_{\kk}W$, $K=\overline{\kk}$, and $R_K=R\otimes_{\kk}K$. For every sufficiently general subspace $L\subseteq [R_K]_1$ of codimension $p+1$, there is a graded surjection $K[x,y]^{(p)}\to R_K/(L)$.
\end{lemma}
\begin{proof}
Fix $0\ne v_1\in E$ and take $v\in E\setminus\kk v_1$. The fraction $r=v/v_1$ sends $(v_1W)R$ into $R$. If this ideal contained a regular sequence $a,b$, then $b(ra)=a(rb)$ would give $ra\in aR$, hence $r\in R_0=\kk$. Thus the nonzero ideal $(v_1W)R$ has height one, since $R$ is Cohen--Macaulay. Choose a height-one prime $\mathfrak p$ containing it. The height hypothesis on $(EW)R$ gives $v_2\in E$ and $w\in W$ with $v_2w\notin\mathfrak p$. In particular, $v_2w\notin v_1W$, so $\dim_{\kk}R_1\geq p+1$.

Choose a basis $u_1,\ldots,u_p$ of $W$ and form the matrix
\[
 M=\begin{pmatrix}v_1u_1&v_1u_2&\cdots&v_1u_p\\v_2u_1&v_2u_2&\cdots&v_2u_p\end{pmatrix}.
\]
A matrix of linear forms is \emph{$1$-generic} over a field if no invertible row and column operations over that field produce a zero entry. For $M$, this means that every nonzero linear combination of its rows has linearly independent entries. We verify this over $\overline{\kk}$. The ratio $t=v_1/v_2=(v_1w)/(v_2w)$ is a nonzero element of $\mathfrak pR_{\mathfrak p}$. If $t$ were algebraic over $\kk$, then $\kk[t]$ would be a field and $t$ would be a unit in $R_{\mathfrak p}$. Thus $t$ is transcendental. For each $\lambda\in\overline{\kk}$, let $g$ be its minimal polynomial over $\kk$. The element $g(t)$ is nonzero and homogeneous of degree zero, and hence a unit in $Q^*(R)$. In $Q^*(R)\otimes_{\kk}\overline{\kk}$ it is divisible by $t-\lambda$, so $t-\lambda$ is a unit there as well. Consequently multiplication by $av_1+bv_2=v_2(at+b)$ is injective on $W\otimes_{\kk}\overline{\kk}$ for all $[a:b]\in\mathbb P^1(\overline{\kk})$. Thus $M$ is $1$-generic over $\overline{\kk}$.

Work now over $K$ and suppress scalar extension from the notation. The map $\mathbb P^1\times\mathbb P(W)\to\mathbb P(R_1)$ sending $([a:b],[w])$ to $[(av_1+bv_2)w]$ is everywhere defined by $1$-genericity. Its image has dimension at most $p$, since $\mathbb P(W)\cong\mathbb P^{p-1}$. Each successive general hyperplane lowers the dimension of the intersection by one, and a general hyperplane misses a finite set of points. Thus $p+1$ general linear functionals on $R_1$ have no common zero on the image. Their common kernel $L$ is a general subspace of codimension $p+1$ and contains no nonzero product $(av_1+bv_2)w$. Hence every nonzero linear combination of the rows of $M$ still has linearly independent entries modulo $L$, so the reduced matrix is $1$-generic.

The polynomial ring on a basis of $R_1/L$ maps onto $R/(L)$. The entries of the reduced matrix are now linear forms in $p+1$ variables. Its $2\times2$ minors vanish in $R/(L)$ because $(v_1u_i)(v_2u_j)=(v_1u_j)(v_2u_i)$ in $R$. If $p\geq2$, these minors generate the ideal of a rational normal curve by \cite[Corollary~6.9~(b)]{E05}. Its coordinate ring $\overline{\kk}[x,y]^{(p)}$ therefore maps onto $R/(L)$. If $p=1$, the polynomial ring on $R_1/L$ itself has two variables.
\end{proof}

We next show that every nonzero graded Artinian algebra of the form $\kk[x,y]^{(p)}/J$ has a unimodal Hilbert function.

\begin{lemma}\label{lem:binary-veronese}
Let $S=\kk[x,y]$, let $p\geq1$, and let $J$ be a homogeneous ideal of $S^{(p)}$ such that $A=S^{(p)}/J$ is nonzero and Artinian. Then $S/JS$ is Artinian and $A=(S/JS)^{(p)}$. Write $H_A(i)=\dim_{\kk}A_i$. Then $H_A$ is unimodal. For each $j\geq0$ such that $H_A(j)<pj+1$, we have $H_A(i+1)\leq H_A(i)$ for all $i\geq j$.
\end{lemma}
\begin{proof}
For each $j\geq0$ we have $[JS]_{pj}=\sum_{i=0}^j J_i S_{p(j-i)}=J_j$, since $S_{p(j-i)}$ is the degree-$(j-i)$ component of $S^{(p)}$ and the term with $i=j$ is $J_j$. Thus $A=C^{(p)}$, where $C=S/JS$. The ring $C$ is finite over $A$, since $S$ is finite over $S^{(p)}$, so $C$ is Artinian.

We may assume that $\kk$ is infinite, since extending the field preserves Hilbert functions. Write $H_C(n)=\dim_{\kk}C_n$ and let $e$ be the least nonnegative integer with $H_C(e)<e+1$. By \cite[Lemma~2.3]{I05}, $H_C(n)=n+1$ for every $0\leq n<e$, and $H_C(n+1)\leq H_C(n)$ for every $n\geq e$. The subsequence $H_A(i)=H_C(pi)$ is therefore unimodal. If $H_A(j)<pj+1$, then $pj\geq e$, so $H_A(i+1)=H_C(pi+p)\leq H_C(pi)=H_A(i)$ for every $i\geq j$.
\end{proof}

To estimate the middle entry, we compare $h_i$ and $h_{s+1-i}$.

\begin{lemma}\label{lem:opposite-entries}
Let $R$ be a $d$-dimensional standard graded Cohen--Macaulay domain over a field $\kk$ with $h$-vector $(1,h_1,\ldots,h_s)$, where $d,s\geq1$. Suppose that $[\tr_R(\omega_R)]_1\ne0$ and that $0\leq h_1-h_s\leq1$. Then
\begin{equation}\label{eq:shifted-h-comparison}
 0\leq h_i-h_{s+1-i}\leq h_1-h_s\qquad\text{for every }i\text{ with }1\leq i\leq\lceil s/2\rceil.
\end{equation}
\end{lemma}
\begin{proof}
Choose a graded fractional ideal $I\cong\omega_R(d-s)$, so that $\indeg I=0$ and $\dim_{\kk} I_0=h_s$. Lemma~\ref{lem:linear-trace-product} gives $I_0[I^{-1}]_1=[\tr_R(\omega_R)]_1\ne0$. Fix $0\ne v\in[I^{-1}]_1$. The ideal $J=vI\cong\omega_R(d-s-1)$ is a proper canonical ideal. Thus $G=R/J$ is Gorenstein of dimension $d-1$ by \cite[Proposition~3.3.18~(b)]{BH98}. Write $G=S/Q$, where $S$ is the polynomial ring on a basis of $G_1$. Since $J_1=vI_0$ has dimension $h_s$, we have $\dim_{\kk}G_1=d+h_1-h_s$. Hence
\[
 c:=\height Q=\dim_{\kk}G_1-\dim G=(d+h_1-h_s)-(d-1)=1+h_1-h_s.
\]
The exact sequence $0\to J\to R\to G\to0$ and graded duality give
\[
 (1-t)h_G(t)=h_R(t)-t^{s+1}h_R(t^{-1}).
\]
For $1\leq i\leq s$, the coefficient of $t^i$ on the right is $h_i-h_{s+1-i}$.

Since $c\le2$, Serre's theorem at the homogeneous maximal ideal of $S$ gives $c$ minimal generators for $Q$ \cite[Corollary~21.20]{E95}. Thus $Q$ is generated by a homogeneous regular sequence. If $h_1=h_s$, then $Q=(f)$ and $(1-t)h_G(t)=1-t^{\deg f}$. Comparing degrees gives $\deg f=s+1$, so all the indicated differences are zero.

If $h_1=h_s+1$, then $Q$ has two generators. Let $a\leq b$ be their degrees. Then
\[
 (1-t)h_G(t)=(1+t+\cdots+t^{a-1})(1-t^b),\qquad a+b=s+2.
\]
In particular, $b>\lceil s/2\rceil$. For each $i$ in the stated range, the coefficient of $t^i$ is $1$ if $i<a$ and $0$ otherwise.
\end{proof}

In a standard graded nearly Gorenstein ring, $[\tr_R(\omega_R)]_1R=\m_R$. Since $\height\m_R=\dim R$, Theorem~\ref{thm:B} follows from the next statement.

\begin{theorem}\label{thm:height-unimodality}
Let $R$ be a standard graded Cohen--Macaulay domain over a field $\kk$, with $h$-vector $(1,h_1,\ldots,h_s)$, where $s\geq1$ and $h_1-h_s\leq1$. Suppose that $\height([\tr_R(\omega_R)]_1R)\geq2$. Then $h_s\leq h_1$, and $h(R)$ is unimodal with $\max_i h_i-h_{\lceil s/2\rceil}\leq1$. If $\operatorname{char}\kk=0$, then the generic Artinian reduction of $R$ has the strong Lefschetz property. If $R$ is level, then $h(R)$ is concave and alternatingly increasing. In particular, in this case it is log-concave and $h_{\lceil s/2\rceil}=\max_i h_i$. If $h_1=h_s$, then $R$ is level and $h_1=\cdots=h_s$.
\end{theorem}
\begin{proof}
If $s=1$, then every linear Artinian reduction has its socle in degree one, and every nonzero linear form is a strong Lefschetz element. The other assertions follow from $h(R)=(1,h_1)$. Assume $s\geq2$ and set $d=\dim R$. Choose a graded fractional ideal $I$ isomorphic to $\omega_R(d-s)$ as a graded $R$-module. Then $\indeg I=0$ and $\dim_{\kk}I_0=h_s$ by \cite[Corollary~4.4.6~(a)]{BH98}. By Lemma~\ref{lem:linear-trace-product}, $I_0[I^{-1}]_1=[\tr_R(\omega_R)]_1$. Since $[\tr_R(\omega_R)]_1\ne0$ by the height assumption, we have $[I^{-1}]_1\ne0$. If $[I^{-1}]_1=\kk v$ for some nonzero $v$, then $[\tr_R(\omega_R)]_1R\subseteq vI$. But $vI$ is a proper canonical ideal, so it has height one by \cite[Proposition~3.3.18~(b)]{BH98}. Thus $\dim_{\kk}[I^{-1}]_1\geq2$.

Apply Lemma~\ref{lem:product-matrix} over $\kk$ with $E=[I^{-1}]_1$ and $W=I_0$. Since $\dim_{\kk}W=h_s$, we obtain a graded surjection $B=\overline{\kk}[x,y]^{(h_s)}\longrightarrow\overline R/(L)$ for every sufficiently general subspace $L\subseteq\overline R_1$ of codimension $h_s+1$, where $\overline R=R\otimes_{\kk}\overline{\kk}$. Since $h_1-h_s\leq1$, we have $\dim L=d+h_1-h_s-1\leq d$. A basis of a general such $L$ therefore extends to a linear system of parameters $\Theta$ of $\overline R$. Thus $B$ maps onto $A=\overline R/\Theta\overline R$, and $\dim_{\overline{\kk}} A_j=h_j$. Lemma~\ref{lem:binary-veronese} gives the unimodality of $h(R)$.

If $\operatorname{char}\kk=0$, Lemma~\ref{lem:binary-veronese} gives $A=C^{(h_s)}$ with $C=\overline{\kk}[x,y]/J$ Artinian. Take a strong Lefschetz element $u\in C_1$ by \cite[Proposition~4.4]{HMNW03}. For every pair $0\le a<b\le s$, multiplication by $(u^{h_s})^{b-a}$ from $A_a$ to $A_b$ is multiplication by $u^{(b-a)h_s}$ from $C_{a h_s}$ to $C_{b h_s}$. Hence $u^{h_s}$ is a strong Lefschetz element of $A$. All linear parameter reductions have Hilbert function $h(R)$, so the minors certifying maximal rank for this choice remain nonzero for generic parameters and a general linear form over their infinite coefficient field. Thus $A_{\mathrm{gen}}$ has the strong Lefschetz property.

Suppose $h_1\leq h_s$. Then $H_A(1)=h_1<h_s+1$. Lemma~\ref{lem:binary-veronese} gives $h_1\geq h_2\geq\cdots\geq h_s$, so $h_1=\cdots=h_s$. If $R$ is Gorenstein, it is level. Otherwise, $\omega_R$ has $h_s=h_1$ minimal generators in its lowest degree. The height assumption gives a nonzero $x\in[\tr_R(\omega_R)]_1$. Since $R$ is a domain, $x$ is a nonzerodivisor, so $R$ is level by \cite[Proposition~4.4~(2)]{MO}.

It remains to consider $h_1=h_s+1$. By Lemma~\ref{lem:opposite-entries}, $0\leq h_i-h_{s+1-i}\leq1$ for every $i$ with $1\leq i\leq\lceil s/2\rceil$. Hence $|h_i-h_{s+1-i}|\leq1$ for every $i$ with $1\leq i\leq s$. Since $h_1=h_s+1>h_0$, we can choose $m\geq1$ with $h_m=\max_i h_i$. The index $\lceil s/2\rceil$ lies between $m$ and $s+1-m$. Unimodality therefore gives $h_{\lceil s/2\rceil}\geq\min\{h_m,h_{s+1-m}\}\geq h_m-1$.

Now suppose that $R$ is level. Then $\overline R$ and its Artinian reduction $A$ are level as well. Every nonzero homogeneous element of $A$ has a nonzero multiple in $A_s$, since a nonzero homogeneous multiple of largest degree lies in the socle. The images of $x^{h_s}$ and $y^{h_s}$ under $B\to A$ lie in $A_1$ by the Veronese grading. For every $2\leq j\leq s$, consider the maps
\[
 A_{j-2}\xrightarrow{\ (y^{h_s},-x^{h_s})\ }A_{j-1}^{\oplus2}\xrightarrow{\ (x^{h_s},y^{h_s})\ }A_j .
\]
Their composite is zero. For every $k\geq2$, each monomial of degree $k h_s$ in $x,y$ is divisible by $x^{h_s}$ or $y^{h_s}$, since a monomial divisible by neither has degree at most $2h_s-2$. Thus $A_k=x^{h_s}A_{k-1}+y^{h_s}A_{k-1}$, so the second map is surjective. If $z\in A_{j-2}$ is in the kernel of the first map, then $x^{h_s}z=y^{h_s}z=0$. As $s-j+2\geq2$, the same equality with $k=s-j+2$ gives $z A_{s-j+2}=0$. By levelness $z=0$, so the first map is injective. Hence $h_{j-2}+h_j\leq2h_{j-1}$ for every $2\leq j\leq s$. Thus $h(R)$ is concave. For every $1\leq i<s$, concavity and the arithmetic--geometric mean inequality give $h_i\geq(h_{i-1}+h_{i+1})/2\geq\sqrt{h_{i-1}h_{i+1}}$. Squaring gives $h_i^2\geq h_{i-1}h_{i+1}$.

It remains to prove $h_i\leq h_{s-i}\leq h_{i+1}$ for every $i$ with $0\leq i<s/2$. The second inequality follows from~\eqref{eq:shifted-h-comparison} at $i+1$. If $i=0$, the first inequality is $1\leq h_s$. Suppose that $1\leq i<s/2$ and $h_i>h_{s-i}$. Then $h_k-h_{k-1}\leq-1$ for some $k$ with $i<k\leq s-i$. By concavity all subsequent differences are at most $-1$, so $h_{s-i+1}\leq h_{s-i}-1\leq h_i-2$. This contradicts~\eqref{eq:shifted-h-comparison} at $i$.
\end{proof}

\begin{example}\label{ex:middle-gap}
Consider the following two projective monomial curves, with homogeneous coordinate rings
\[
\begin{aligned}
 R&=\QQ[u^{21},u^{17}v^4,u^{12}v^9,u^9v^{12},u^8v^{13},v^{21}],\\
 S&=\QQ[u^{27},u^{22}v^5,u^{15}v^{12},u^{12}v^{15},u^{10}v^{17},v^{27}].
\end{aligned}
\]
In each ring the displayed generators have degree one. The ring $R$ is the type-four example constructed in \cite[proof of Theorem~3.11, case $d=4$]{M24}. In particular, $R$ is nearly Gorenstein and is not level. Its $h$-vector is $(1,4,7,6,3)$, with its maximum at the middle entry.

For $S$, one can check with \texttt{Macaulay2}~\cite{M2} that $S$ is Cohen--Macaulay, $\tr_S(\omega_S)=\m_S$, and $h(S)=(1,4,7,6,6,3)$. Thus $S$ is nearly Gorenstein and $h_1-h_5=1$. The middle entry specified in Theorem~\ref{thm:height-unimodality} is $h_3=6$, while $h_2=7$ is the maximum. This gives equality in the bound $\max_i h_i-h_{\lceil s/2\rceil}\le1$. Moreover, $h_3^2=36<42=h_2h_4$, so $h(S)$ is not log-concave. Thus $S$ is not level by Theorem~\ref{thm:height-unimodality}.
\end{example}

\begin{remark}\label{rem:yanagawa}
Over an algebraically closed field of characteristic zero, a standard graded Cohen--Macaulay domain with $s\geq2$ and $h_1=\cdots=h_{s-1}\geq h_s\geq2$ is level \cite[Corollary~3.11]{Y95}, \cite[Lemma~5~(2) and Proposition~6]{Y97}.
\end{remark}

\section*{Acknowledgments}
The first author was partially supported by JSPS KAKENHI Grant Number JP25K07100. The second author was supported by JSPS KAKENHI Grant Number JP25KJ1744 and thanks Professor Sugiyama for the opportunity that led to this collaboration. The authors used ChatGPT to examine proofs, perform computations, check references, and revise the text; they take full responsibility for all mathematical content.\par

\begin{thebibliography}{AKOPP26}
\setlength{\baselineskip}{10.35pt}
\setlength{\itemsep}{0.6pt}

\bibitem[AKOPP26]{AKOPP26} Karim Alexander Adiprasito, Eric Katz, Ryoshun Oba, Stavros Argyrios Papadakis, and Vasiliki Petrotou, Frobenius identities for the volume map on Cohen--Macaulay rings, preprint, arXiv: 2605.02479v1, 2026.

\bibitem[APP25]{APP25} Karim Alexander Adiprasito, Stavros Argyrios Papadakis, and Vasiliki Petrotou, Lattice polytopes and semigroup algebras: Generic Lefschetz properties and Parseval--Rayleigh identities, preprint, arXiv: 2509.14152, 2025.

\bibitem[APPS22]{APPS22} Karim Alexander Adiprasito, Stavros Argyrios Papadakis, Vasiliki Petrotou, and Johanna Kristina Steinmeyer, Beyond positivity in Ehrhart theory, preprint, arXiv: 2210.10734, 2022.

\bibitem[AY21]{AY21} Karim Adiprasito and Geva Yashfe, The partition complex: an invitation to combinatorial commutative algebra, in Surveys in Combinatorics 2021, London Mathematical Society Lecture Note Series, vol.~470, Cambridge University Press, Cambridge, 2021, 1--42.

\bibitem[A26]{Ath26} Christos A. Athanasiadis, Triangulations of simplicial complexes and theta polynomials, Tohoku Mathematical Journal (2) 78 (2026), no.~3, 423--441.

\bibitem[BH98]{BH98} Winfried Bruns and J\"urgen Herzog, Cohen--Macaulay rings, revised edition, Cambridge Studies in Advanced Mathematics, vol. 39, Cambridge University Press, Cambridge, 1998.

\bibitem[BK05]{BK05} Michel Brion and Shrawan Kumar, Frobenius splitting methods in geometry and representation theory, Progress in Mathematics, vol. 231, Birkh\"auser, Boston, 2005.

\bibitem[BR07]{BR07} Winfried Bruns and Tim R\"omer, $h$-Vectors of Gorenstein polytopes, Journal of Combinatorial Theory, Series A 114 (2007), no. 1, 65--76.

\bibitem[CSV25]{CSV25} Aldo Conca, Anurag K. Singh, and Matteo Varbaro, Invariant rings of the special orthogonal group have nonunimodal $h$-vectors, Journal of Algebra and Its Applications 24 (2025), nos. 13--14, 2540001, 10 pp.

\bibitem[E95]{E95} David Eisenbud, Commutative algebra with a view toward algebraic geometry, Graduate Texts in Mathematics, vol.~150, Springer-Verlag, New York, 1995.

\bibitem[E05]{E05} David Eisenbud, The geometry of syzygies: A second course in commutative algebra and algebraic geometry, Graduate Texts in Mathematics, vol. 229, Springer, New York, 2005.

\bibitem[EY26]{EY26} Naoki Endo and Ken-ichi Yoshida, Nearly Gorenstein blow-up algebras over two-dimensional regular local rings, preprint, arXiv:2607.24019, 2026.

\bibitem[F26]{F26} Luis Ferroni, Unimodality shenanigans in Ehrhart theory, preprint, arXiv:2609.10513v2, 2026.

\bibitem[GW78]{GW78} Shiro Goto and Keiichi Watanabe, On graded rings, I, Journal of the Mathematical Society of Japan 30 (1978), no.~2, 179--213.

\bibitem[M2]{M2} Daniel R. Grayson and Michael E. Stillman, \texttt{Macaulay2}, a software system for research in algebraic geometry, available at \url{https://macaulay2.com}.

\bibitem[HKMM23]{HKMM23} Thomas Hall, Max K\"olbl, Koji Matsushita, and Sora Miyashita, Nearly Gorenstein polytopes, Electronic Journal of Combinatorics 30 (2023), no. 4, Paper No. P4.42, 21 pp.

\bibitem[HMNW03]{HMNW03} Tadahito Harima, Juan C. Migliore, Uwe Nagel, and Junzo Watanabe, The weak and strong Lefschetz properties for Artinian $K$-algebras, Journal of Algebra 262 (2003), no.~1, 99--126.

\bibitem[HHS19]{HHS19} J\"urgen Herzog, Takayuki Hibi, and Dumitru I. Stamate, The trace of the canonical module, Israel Journal of Mathematics 233 (2019), no. 1, 133--165.

\bibitem[HMP19]{HMP19} J\"urgen Herzog, Fatemeh Mohammadi, and Janet Page, Measuring the non-Gorenstein locus of Hibi rings and normal affine semigroup rings, Journal of Algebra 540 (2019), 78--99.

\bibitem[H89]{H89} Takayuki Hibi, Flawless $O$-sequences and Hilbert functions of Cohen--Macaulay integral domains, Journal of Pure and Applied Algebra 60 (1989), no. 3, 245--251.

\bibitem[H92]{H92} Takayuki Hibi, Algebraic combinatorics on convex polytopes, Carslaw Publications, Glebe, 1992.

\bibitem[Hi19]{Hi19} Akihiro Higashitani, Unimodality of $\delta$-vectors of lattice polytopes and two related properties, European Journal of Mathematics 5 (2019), no.~2, 333--355.

\bibitem[Hoc72]{Hoc72} Melvin Hochster, Rings of invariants of tori, Cohen--Macaulay rings generated by monomials, and polytopes, Annals of Mathematics (2) 96 (1972), no. 2, 318--337.

\bibitem[I05]{I05} Anthony Iarrobino, Betti strata of height two ideals, Journal of Algebra 285 (2005), no.~2, 835--855.

\bibitem[KKP+21]{KKP21} Zhibek Kadyrsizova, Jennifer Kenkel, Janet Page, Jyoti Singh, Karen E. Smith, Adela Vraciu, and Emily E. Witt, Cubic surfaces of characteristic two, Transactions of the American Mathematical Society 374 (2021), no.~9, 6251--6267.

\bibitem[KLS26]{KLS26} Kalle Karu, Matt Larson, and Alan Stapledon, Differential operators, anisotropy, and simplicial spheres, International Mathematics Research Notices 2026 (2026), no.~13, rnag147.

\bibitem[M24]{M24} Sora Miyashita, Levelness versus nearly Gorensteinness of homogeneous rings, Journal of Pure and Applied Algebra 228 (2024), no.~4, 107553.

\bibitem[M26a]{M26} Sora Miyashita, A linear generalization of the nearly Gorenstein property, with applications to Veronese subalgebras, Mathematische Nachrichten 299 (2026), no. 9, 2436--2451.

\bibitem[M26b]{MPG26} Sora Miyashita, When do pseudo-Gorenstein rings become Gorenstein?, Bulletin of the London Mathematical Society 58 (2026), no.~1, e70186.

\bibitem[MO]{MO} Sora Miyashita and Taiga Ozaki, Higher-dimensional Teter rings via the canonical trace, to appear in Journal of the Mathematical Society of Japan, arXiv:2512.06761.

\bibitem[MP05]{MP05} Mircea Musta\c{t}\u{a} and Sam Payne, Ehrhart polynomials and stringy Betti numbers, Mathematische Annalen 333 (2005), no. 4, 787--795.

\bibitem[OH06]{OH06} Hidefumi Ohsugi and Takayuki Hibi, Special simplices and Gorenstein toric rings, Journal of Combinatorial Theory, Series A 113 (2006), no. 4, 718--725.

\bibitem[Pa19]{Pa19} Janet Page, The Frobenius complexity of Hibi rings, Journal of Pure and Applied Algebra 223 (2019), no. 2, 580--604.

\bibitem[PP20]{PP20} Stavros Argyrios Papadakis and Vasiliki Petrotou, The characteristic $2$ anisotropicity of simplicial spheres, preprint, arXiv:2012.09815, 2020.

\bibitem[P08]{P08} Sam Payne, Ehrhart series and lattice triangulations, Discrete \& Computational Geometry 40 (2008), no. 3, 365--376.

\bibitem[Sch22]{Sch22} Stefan Schr\"oer, Elliptic curves over the rational numbers with semi-abelian reduction and two-division points, Journal of Number Theory 231 (2022), 80--101.

\bibitem[Si00]{Si00} Anurag K. Singh, Multi-symbolic Rees algebras and strong $F$-regularity, Mathematische Zeitschrift 235 (2000), 335--344.

\bibitem[S96]{S96} Richard P. Stanley, Combinatorics and commutative algebra, second edition, Progress in Mathematics, vol.~41, Birkh\"auser Boston, 1996.

\bibitem[SVL13]{SVL13} Jan Schepers and Leen Van Langenhoven, Unimodality questions for integrally closed lattice polytopes, Annals of Combinatorics 17 (2013), 571--589.

\bibitem[S78]{S78} Richard P. Stanley, Hilbert functions of graded algebras, Advances in Mathematics 28 (1978), no. 1, 57--83.

\bibitem[S89]{S89} Richard P. Stanley, Log-concave and unimodal sequences in algebra, combinatorics, and geometry, Annals of the New York Academy of Sciences 576 (1989), 500--535.

\bibitem[Y95]{Y95} Kohji Yanagawa, Castelnuovo's lemma and $h$-vectors of Cohen--Macaulay homogeneous domains, Journal of Pure and Applied Algebra 105 (1995), no. 1, 107--116.

\bibitem[Y97]{Y97} Kohji Yanagawa, A note on the $h$-vector of a Cohen--Macaulay homogeneous domain, Japanese Journal of Mathematics (N.S.) 23 (1997), no. 1, 179--186.

\end{thebibliography}

\end{document}



